\begin{center}
\textsc{Gauss} an \textsc{Bessel}.\quad Göttingen, 12.~März 1826.
\end{center}

\par\noindent\hspace*{\parindent}.\ .\ .\ .\ .\ .\quad Den größten Teil des vergangenen Sommers habe ich mit meinen Messungen in den Gebieten von Bremen und Oldenburg verbracht.\quad.\ .\ .\ .\ .\ .

\par\noindent\hspace*{\parindent}Was meine Messungen betrifft, so habe ich deren trigonometrischen Teil zumindest dem Buchstaben nach vollendet. Meine Seite Varel-Jever schließt sich an die \textsc{Krayenhoff}schen Dreiecke an.
Ob die Messungen aber wirklich abgeschlossen sind, weiß ich selbst noch nicht.
Meine Winkelmessungen in Jever ergeben \textls[120]{ganz enorme} Unterschiede gegenüber den \textsc{Krayenhoff}schen, die bis zu \(15^{\prime\prime}\) betragen.
Ihre Ursache
\par\clearpage
\noindent{\small\textbf{ZUR HANNOVERSCHEN TRIANGULATION.}}\hfill361\par\medskip
\noindent kann ich nicht mit Sicherheit angeben.
Für meine eigenen Winkel bürge ich bis auf eine Bogensekunde.
Die Unterschiede ließen sich erklären, wenn ich annehmen dürfte, dass \textsc{Krayenhoffs} Stationszentrum von meinem (ich habe ein Stockwerk tiefer beobachtet) einen Meter entfernt liegt.
Die Auskünfte, die ich von \textsc{Krayenhoff} erhalten habe, sind hinsichtlich der Zentrierung unbefriedigend. Im Allgemeinen geht daraus aber hervor, dass seine Winkelmessungen in jener Gegend bei Weitem nicht die Genauigkeit meiner Messungen haben.
Er hat in Ostfriesland ein schlechteres Instrument verwendet. Von vielen Winkelreihen hat er immer nur diejenigen beibehalten, die am besten zu passen schienen (ohne anzugeben, um wie viel die anderen abwichen). Selbst unter den in Jever beibehaltenen finden sich Unterschiede von \(4^{\prime\prime}\).
Unter diesen Umständen scheint es mir nicht ratsam, meine Messungen in Ostfriesland weiter auszudehnen.
Seine weiter südlich vorgenommenen Messungen sind besser, und eine Verbindung meiner Dreiecke über das Osnabrücker Gebiet nach Bentheim könnte zweifellos zuverlässigere Ergebnisse liefern.
Der schwankende Zustand meiner Gesundheit hat mir aber bisher den Mut genommen, solche Arbeiten vorzuschlagen, die noch einen ein- oder anderthalbjährigen Feldeinsatz erfordern würden. In diesem Sommer wird sich in dieser Hinsicht kaum etwas Wesentliches tun lassen.\quad.\ .\ .\ .\ .\ .

